\subsection{Induced and Coinduced Representations}

Let H be a subgroup of the group G.

If \((V,\pi)\) is a linear representation of G, then the restriction of \(\pi\) to H is a linear representation of H in V; we denote it by \(\operatorname{Res}_{\mathrm{H}}^{\mathrm{G}}(\pi)\) . The K{[}H{]}-module associated with \(\operatorname{Res}_{\mathrm{H}}^{\mathrm{G}}(\pi)\) is simply the module deduced from the K{[}G{]}-module V by restriction of scalars (II, §1, No.~13, p.~221). If V is a finite-dimensional free K-module, then the character of \(\operatorname{Res}_{\mathrm{H}}^{\mathrm{G}}(\pi)\) is the restriction of the character of \(\pi\) to H.

Let \((\mathrm{M},\sigma)\) be a representation of H.

We view K{[}G{]} as a (K{[}G{]}, K{[}H{]})-bimodule and M as a K{[}H{]}-module. The K{[}G{]}-module \(\mathcal{T}(M) = K[G] \otimes_{K[H]} M\) (VIII, p.~58) defines a linear representation of G, denoted by \(\operatorname{Ind}_{\mathrm{H}}^{\mathrm{G}}(\sigma)\) and called the representation of G induced by \(\sigma\) . If \((\mathrm{V}, \pi)\) is a linear representation of G, then the K{[}H{]}-module \(\mathcal{H}(\mathrm{V}) = \mathrm{Hom}_{\mathrm{K[G]}}(\mathrm{K[G]},\mathrm{V})\) can be identified with the K{[}H{]}-module corresponding to the representation \(\mathrm{Res}_{\mathrm{H}}^{\mathrm{G}}(\pi)\) . Consequently, the adjunction morphism (VIII, p.~59) gives a K-module isomorphism, called canonical, from \(\mathrm{Hom}_{\mathrm{H}}(\sigma ,\mathrm{Res}_{\mathrm{H}}^{\mathrm{G}}(\pi))\) to \(\mathrm{Hom}_{\mathrm{G}}(\mathrm{Ind}_{\mathrm{H}}^{\mathrm{G}}(\sigma),\pi)\) (``Frobenius reciprocity'').

We view K{[}G{]} as a (K{[}H{]}, K{[}G{]})-bimodule. The K{[}G{]}-module \(\mathcal{H}(M) = \operatorname{Hom}_{\mathrm{K}[\mathrm{H}]}(\mathrm{K}[\mathrm{G}], \mathrm{M})\) defines a representation of G, denoted by \(\operatorname{Coind}_{\mathrm{H}}^{\mathrm{G}}(\sigma)\) and called the representation of G coinduced by \(\sigma\) . If \((\mathrm{V}, \pi)\) is a linear representation of G, then the K{[}H{]}-module \(\mathcal{T}(\mathrm{V}) = \mathrm{K}[\mathrm{G}] \otimes_{\mathrm{K}[\mathrm{G}]} \mathrm{V}\) can be identified with the K{[}H{]}-module corresponding to the representation \(\operatorname{Res}_{\mathrm{H}}^{\mathrm{G}}(\pi)\) . Consequently, the adjunction morphism (loc. cit.) gives a K-module isomorphism, called canonical, from \(\operatorname{Hom}_{\mathrm{H}}(\operatorname{Res}_{\mathrm{H}}^{\mathrm{G}}(\pi), \sigma)\) to \(\operatorname{Hom}_{\mathrm{G}}(\pi, \operatorname{Coind}_{\mathrm{H}}^{\mathrm{G}}(\sigma))\) .

Let \(\varepsilon : \mathrm{K}[\mathrm{G}] \to \mathrm{K}[\mathrm{H}]\) be the K-module homomorphism characterized by the relations \(\varepsilon(h) = h\) if \(h \in \mathrm{H}\) and \(\varepsilon(g) = 0\) if \(g \in \mathrm{G} - \mathrm{H}\) . The mapping \(\varepsilon\) is a homomorphism of \((\mathrm{K}[\mathrm{H}], \mathrm{K}[\mathrm{H}])\) -bimodules. Let \((\mathrm{M}, \sigma)\) be a linear representation of H. The mapping \(v \mapsto v \circ \varepsilon\) from \(\mathrm{Hom}_{\mathrm{K}[\mathrm{H}]}(\mathrm{K}[\mathrm{H}], \mathrm{M})\) to \(\mathrm{Hom}_{\mathrm{K}[\mathrm{H}]}(\mathrm{K}[\mathrm{G}], \mathrm{M})\) is a homomorphism of \(\mathrm{K}[\mathrm{H}]\) -modules. By identifying M with \(\mathrm{Hom}_{\mathrm{K}[\mathrm{H}]}(\mathrm{K}[\mathrm{H}], \mathrm{M})\) , we obtain a homomorphism of \(\mathrm{K}[\mathrm{H}]\) -modules from M to \(\mathrm{Res}_{\mathrm{H}}^{\mathrm{G}}(\mathrm{Coind}_{\mathrm{H}}^{\mathrm{G}}(\sigma))\) . Frobenius reciprocity sends it to a homomorphism \(\iota\) of \(\mathrm{K}[\mathrm{G}]\) -modules from \(\mathrm{Ind}_{\mathrm{H}}^{\mathrm{G}}(\sigma)\) to \(\mathrm{Coind}_{\mathrm{H}}^{\mathrm{G}}(\sigma)\) characterized by the relation

\[
\iota (g \otimes m) (g ^ {\prime}) = \varepsilon (g ^ {\prime} g) m
\]

for \(g, g' \in G\) and \(m \in M\) . This homomorphism is called canonical.

PROPOSITION 2. --- Let H be a subgroup of G, and let \((\mathrm{M},\sigma)\) be a linear representation of H. The canonical homomorphism \(\iota:\operatorname{Ind}_{\mathrm{H}}^{\mathrm{G}}(\sigma)\to\operatorname{Coind}_{\mathrm{H}}^{\mathrm{G}}(\sigma)\) is injective. If the subgroup H has finite index in G, then \(\iota\) is an isomorphism of K{[}G{]}-modules.

Let \(S \subset G\) be a system of representatives of \(G/H\) . The family \((s)_{s \in S}\) is a basis of the right \(K[H]\) -module \(K[G]\) . It follows that the mapping \(M^{(S)} \to \operatorname{Ind}_{H}^{G}(\sigma)\) defined by \((m_s)_{s \in S} \mapsto \sum_{s \in S} s \otimes m_s\) is an isomorphism of \(K\) -modules. For all \(s, s' \in S\) and every \(m \in M\) , we have the relations

\[
\iota (s ^ {\prime} \otimes m) (s ^ {- 1}) = \left\{ \begin{array}{l} m \text {   if   } s = s ^ {\prime}, \\ 0 \text {   otherwise }. \end{array} \right.
\]

It follows that \(\iota'\) is injective.

Suppose that H has finite index in G. The set S is then finite; let \(\rho\) be the mapping from \(\operatorname{Coind}_{\mathrm{H}}^{\mathrm{G}}(\sigma)\) to \(\operatorname{Ind}_{\mathrm{H}}^{\mathrm{G}}(\sigma)\) given by \(u \mapsto \sum_{s \in S} s \otimes u(s^{-1})\) . It satisfies \((\iota \circ \rho(u))(s^{-1}) = u(s^{-1})\) for \(u \in \operatorname{Coind}_{\mathrm{H}}^{\mathrm{G}}(\sigma)\) and \(s \in S\) . Since the family \((s^{-1})_{s\in\mathrm{S}}\) is a basis of the left K{[}H{]}-module K{[}G{]}, the mapping \(\iota\circ\rho\) is the identity mapping. Consequently, \(\iota\) is bijective with inverse bijection \(\rho\) .

Let H be a subgroup of G of finite index. Let u be a central function on H; denote by \(u^{0}\) the function on G that extends u and that is zero at every point of G---H. Let S be a system of representatives of G/H. For \(g \in G\) , set

\[
\operatorname{Ind} _ {\mathrm{H}} ^ {\mathrm{G}} (u) (g) = \sum_ {s \in \mathrm{S}} u ^ {0} (s ^ {- 1} g s).\tag{4}
\]

Note that for all \(x, g \in G\) and every \(h \in H\) , we have \(u^{0}((xh)^{-1}gxh) = u^{0}(x^{-1}gx)\) . It follows that \(\operatorname{Ind}_{\mathrm{H}}^{\mathrm{G}}(u)\) is a central function on G that does not depend on the choice of S. We thus define a K-linear mapping \(Ind_{H}^{G}\) from \(\mathcal{Z}_{\mathrm{K}}(\mathrm{H})\) to \(\mathcal{Z}_{\mathrm{K}}(\mathrm{G})\) .

PROPOSITION 3. --- Let H be a subgroup of G of finite index. Let \((M,\sigma)\) be a linear representation of H. Suppose that M is a finite-dimensional free K-module. Denote by \((V,\pi)\) the representation of G induced by \(\sigma\) . Then the K-module V is free and finite-dimensional, and

\[
\chi_ {\pi} = \mathrm{Ind} _ {\mathrm{H}} ^ {\mathrm{G}} (\chi_ {\sigma}).\tag{5}
\]

Let S be a system of representatives of G/H. As we saw above, the linear mapping \(\mathrm{M}^{\mathrm{S}}\to\mathrm{Ind}_{\mathrm{H}}^{\mathrm{G}}(\mathrm{M})\) given by \((m_{s})_{s\in\mathrm{S}}\mapsto\sum_{s\in\mathrm{S}}s\otimes m_{s}\) is a K-module isomorphism. In particular, V is a finite-dimensional free K-module. For any \(g\in G\) , denote by \(M_{g}\) the image of M by the mapping \(m\mapsto g\otimes m\) . For all \(g,g'\in G\) , we have \(M_{g}=M_{g'}\) if and only if \(gH=g'H\) , and V is the direct sum of the submodules \(M_{s}\) for \(s\in S\) . For all \(g,g'\in G\) , we also have \(\pi(g)\mathrm{M}_{g'}=\mathrm{M}_{gg'}\) . For any \(g\in G\) , denote by \(S_{g}\) the set of elements s of S such that \(s^{-1}gs\in H\) . For all \(g,g'\in G\) such that \(g'^{-1}gg'\in H\) , the automorphism \(\pi(g)\) of V induces an automorphism \(\pi(g)_{g'}\) of \(M_{g'}\) , and we have

\[
\operatorname{Tr} (\pi (g)) = \sum_ {s \in \mathrm{S} _ {g}} \operatorname{Tr} (\pi (g) _ {s}) = \sum_ {s \in \mathrm{S} _ {g}} \operatorname{Tr} (\sigma (s ^ {- 1} g s)).
\]

The last assertion of the proposition follows.

